% MODELO DE SUBMISSÃO DE RESUMO PARA APRESENTACAO DE TRABALHO
% E DE PÔSTER NO XVIV ENCONTRO CIENTÍFICO DE PÓS-GRADUAÇÃO

% BASEADO NOS TEMPLATES OFICIAIS DA SBMAC PARA O CNMAC 2016.
% PARA MAIS INSTRUÇÕES SOBRE FIGURAS, TABELAS E EQUAÇÕES,
% BAIXE O ARQUIVO DISPONIVEL EM
% http://cnmac.org.br/br/node/264

\documentclass[a4paper, 12pt]{article}

\usepackage[brazil]{babel}      % para texto em Portugues
%\usepackage[english]{babel}    % para texto em Ingles

\usepackage[utf8]{inputenc}   % para acentuação em Português com o uso do Unicode, 

%% POR FAVOR, EVITE MUDAR A CONFIGURAÇÃO DESTE MODELO.
%% ADICIONE APENAS OS PACOTES ESSENCIAIS PARA SEU TRABALHO.

\usepackage[centertags]{amsmath}
\usepackage{amssymb,amsfonts,amsthm}
\usepackage{indentfirst}
\usepackage{array}
\usepackage{float}
\usepackage[usenames,dvipsnames,svgnames,table]{xcolor}
\usepackage{hyperref}
\usepackage{enumerate}
\usepackage{graphicx}
\usepackage{epstopdf}
\usepackage{graphics}
\usepackage{subfigure}
\usepackage[margin=1.5cm,footskip=1.5cm,includefoot]{geometry} 
\usepackage{listings}
\usepackage{setspace}
\usepackage{multirow}
\usepackage{lipsum}

\setstretch{1.1}

\begin{document}
	%
	%
	%********************************************************
	% ************ CABEÇALHO - NAO MODIFICAR ****************	
	%********************************************************
	%
	%
	\begin{center}
		\sc\large\textbf{XIV Encontro Científico de Pós-Graduandos do IMECC}\\
		\vspace{-1em}
		
oindent \hrulefill\\	
		\vspace{-1.1em}
		
oindent \hrulefill\\
	\end{center}
	%
	%
	%********************************************************
	%****************** TÍTULO DO TRABALHO ******************	
	%********************************************************	
	%
	%
	\begin{center}
		\Large{Demonstratio nova theorematis omnem functionem algebraicam rationalem integram unius variabilis in factores}
	\end{center}
	%
	%
	%********************************************************
	%*************** INFORMAÇÕES DOS AUTORES ****************	
	%********************************************************	
	%
	%
	% LEMBRE-SE DE SUBLINHAR O NOME DE QUEM VAI REALIZAR A
	% APRESENTAÇÃO
	%
	\author{
		{\large  Nome da Autora }\thanks{autora1@email}\\
		{\small Filiação da autora} \\
		{\large Nome do Autor }\thanks{autor2@email} \\
		{\small Filiação do autor } \\
	}
	
	% PRIMEIRO AUTOR
	
oindent {\small 
	\underline{Johann C. F. Gauss}\footnote{\texttt{dekan@math-cs.uni-goettingen.de}},
	%
	% SEGUNDO AUTOR	
	Isaac Newton\footnote{\texttt{newton@cam.ac.uk}},
	%
	% TERCEIRO AUTOR
	Leonhard Paul Euler\footnote{\texttt{euler@unibas.ch}}}\\
	%
	% ORIGEM DO PRIMEIRO AUTOR
	{\footnotesize
	$^{1}$Departamento de Matemática, University of G\"{o}ttingen, G\"{o}ttingen, Alemanha\\ 
	%
	% ORIGEM DO SEGUNDO AUTOR
	$^{2}$Departamento de Matemática, University of Cambridge, Cambridge, Reino Unido\\
	%
	% ORIGEM DO TERCEIRO AUTOR
	$^{3}$Departamento de Matemática e Informática, University of Basel, Basileia, Suíça}
	%
	%
	%********************************************************
	%**************** DESCRIÇÃO DO TRABALHO *****************	
	%********************************************************	
	%
	% A CRIAÇÃO DE SESSÕES É LIVRE, DESDE QUE RESPEITADO O
	% LIMITE DE DUAS (2) PÁGINAS, INCLUINDO AS REFERÊNCIAS.
	%
	\begin{abstract}
	\lipsum[1]
	
	
oindent {\bf Palavras-chave}:   % De 2 a 4 palavras
	% separadas por vírgulas.
	\end{abstract}
	\section{Introdução}   % Os títulos aqui colocados são
	% apenas sugestões.
	\lipsum[2]
	\section{Desenvolvimento}
	\lipsum[3-6]	
	\section{Considerações Finais}
	\lipsum[7]
	%
	%
	%********************************************************
	%********************* REFERÊNCIAS **********************	
	%********************************************************		
	%
	%
	% POR FAVOR, MANTENHA A BIBLIOGRAFIA NO PADRÃO RECOMENDADO.
	% COLOQUE APENAS AS PRINCIPAIS REFERÊNCIAS DO TRABALHO,
	% SEGUNDO A ORDEM ALFABÉTICA DO SOBRENOME DO PRIMEIRO AUTOR.
	%
	\begin{thebibliography}{00} \footnotesize
		\bibitem{Boldrini}
		J. L. Boldrini, S. I. R. Costa, V. R. Ribeiro, and H. G. Wetzler. {\it Álgebra Linear e Aplicações}. Harper-Row, São Paulo, 1987.
		\bibitem{Cuminato}
		J. A. Cuminato and V. Ruas, Unification of distance inequalities for linear variational problems, {\it Comp. Appl. Math.}, 2014. DOI: 10.1007/s40314-014-0163-6.
		\bibitem{Santos}
		I. L. D. Santos and  G. N. Silva. Uma classe de problemas de controle ótimo em escalas temporais, {\it Proceeding Series of the Brazilian Society of Computational and Applied Mathematics}, volume 1, 2013. DOI: 10.5540/03.2013.001.01.0177.
	\end{thebibliography}
\end{document}